\documentclass[11pt, a4paper]{article}
\usepackage{amsmath, amssymb, amsthm}
\usepackage{geometry}
\usepackage{booktabs}
\usepackage{xcolor}
\usepackage{microtype}
\usepackage{hyperref}

\geometry{margin=2.8cm}

\title{\textbf{SierpSphere Grammar Notation}\\
       \large A compact symbolic language for fractal SDF grammars}
\author{}
\date{}

\newcommand{\G}[1]{\mathsf{#1}}          % group symbol
\newcommand{\prim}[1]{\mathrm{#1}}       % primitive symbol
\newcommand{\sub}{\ominus}               % subtract
\newcommand{\add}{\oplus}                % add
\newcommand{\intr}{\cap}                 % intersect
\newcommand{\step}[4]{#1_{\prim{#2}}^{#3}{}_{#4}}   % op_P^sigma_rho
\newcommand{\stepd}[5]{#1_{\prim{#2}}^{#3}({#4})_{#5}} % op_P^sigma(delta)_rho
\newcommand{\gram}[2]{\G{#1}[\prim{#2}]}  % group[seed]

\begin{document}
\maketitle
\thispagestyle{empty}

% ── 1. Motivation ──────────────────────────────────────────────────────────
\section*{Motivation}

A SierpSphere grammar is fully specified by three things: a \emph{symmetry group},
a \emph{seed primitive}, and an ordered list of \emph{iteration steps}.
The JSON representation is verbose; this notation condenses a grammar to a
single line amenable to comparison, mutation logging, and typeset publication.

% ── 2. Alphabet ────────────────────────────────────────────────────────────
\section*{Alphabet}

\paragraph{Symmetry groups} (Schoenflies notation)
\[
  \G{T_d} \quad \G{O_h} \quad \G{I_h}
\]
corresponding to \textit{tetrahedral}, \textit{octahedral}, and
\textit{icosahedral} point groups respectively.

\paragraph{Primitives}
\[
  \prim{S} = \text{sphere} \qquad
  \prim{C} = \text{cube}   \qquad
  \prim{O} = \text{octahedron}
\]

\paragraph{Operations}
\[
  \sub = \text{smooth-subtract} \qquad
  \add = \text{smooth-union}    \qquad
  \intr = \text{smooth-intersect}
\]

% ── 3. Step notation ───────────────────────────────────────────────────────
\section*{Step notation}

A single iteration step carrying parameters
$(\sigma, \delta, \rho) \in (0,1)\times\mathbb{R}_{>0}\times\mathbb{R}_{\ge 0}$
is written:

\[
  \boxed{\; \mathrm{op}_{\prim{P}}^{\,\sigma}(\delta)_{\rho} \;}
\]

\begin{center}
\begin{tabular}{lll}
\toprule
Symbol & Parameter & Default (omit when) \\
\midrule
$\sigma$ & \texttt{scale\_factor} — child size relative to parent & — (always shown) \\
$\delta$ & \texttt{distance\_factor} — placement along symmetry axis & 1 (omit when $\delta=1$) \\
$\rho$   & \texttt{smooth\_radius} — Boolean blending width & 0 (omit when $\rho=0$) \\
\bottomrule
\end{tabular}
\end{center}

\medskip
\noindent Examples of a single step:
\[
  \sub_{\prim{S}}^{0.5}{}_{\!0.02}
  \qquad
  \add_{\prim{C}}^{0.3}(1.2)_{\!0.05}
  \qquad
  \intr_{\prim{O}}^{0.45}
\]

% ── 4. Full grammar formula ────────────────────────────────────────────────
\section*{Full grammar formula}

\[
  \boxed{\;
    \Gamma \;=\; \G{G}[\prim{P_0}] \;:\; s_1 \cdot s_2 \cdots s_n
  \;}
\]

where $\G{G}$ is the symmetry group, $\prim{P_0}$ the seed primitive, and
$s_1, \ldots, s_n$ are steps applied in order.
Steps act on \emph{new} nodes by default; a trailing ${}^{\!\star}$
marks \texttt{apply\_to: all}.

% ── 5. Named grammars ──────────────────────────────────────────────────────
\section*{Named grammars}

\newcommand{\half}{\tfrac{1}{2}}

\paragraph{\texttt{sierpinski\_classic}}
Classic tetrahedral Sierpi\'nski gasket from a sphere seed,
alternating subtract / union / subtract with decreasing smoothness.
\[
  \Gamma_{\mathrm{classic}}
  = \G{T_d}[\prim{S}]
  : \sub_{\prim{S}}^{\half}{}_{\!0.02}
  \cdot \add_{\prim{S}}^{\half}{}_{\!0.01}
  \cdot \sub_{\prim{S}}^{\half}{}_{\!0.005}
\]

\paragraph{\texttt{sierpinski\_octahedron}}
Same alternating pattern, octahedron seed and primitives under $\G{T_d}$.
\[
  \Gamma_{\mathrm{oct}}
  = \G{T_d}[\prim{O}]
  : \sub_{\prim{O}}^{\half}{}_{\!0.02}
  \cdot \add_{\prim{O}}^{\half}{}_{\!0.01}
  \cdot \sub_{\prim{O}}^{\half}{}_{\!0.005}
\]

\paragraph{\texttt{sierpinski\_cube}}
Pure additive cube fractal — each step unions scaled cubes
along tetrahedral axes.
\[
  \Gamma_{\mathrm{cube}}
  = \G{T_d}[\prim{C}]
  : \add_{\prim{C}}^{\half}{}_{\!0.02}
  \cdot \add_{\prim{C}}^{\half}{}_{\!0.01}
  \cdot \add_{\prim{C}}^{\half}{}_{\!0.005}
\]

\paragraph{\texttt{sierpinski\_void}}
Four-step icosahedral grammar producing high-genus void structure.
\[
  \Gamma_{\mathrm{void}}
  = \G{I_h}[\prim{S}]
  : \sub_{\prim{S}}^{\half}{}_{\!0.02}
  \cdot \add_{\prim{S}}^{\half}{}_{\!0.01}
  \cdot \sub_{\prim{S}}^{\half}{}_{\!0.005}
  \cdot \add_{\prim{S}}^{\half}{}_{\!0.01}
\]

% ── 6. Mutation shorthand ──────────────────────────────────────────────────
\section*{Mutation shorthand}

During evolution, a grammar can be reported as a delta from a named parent:
\[
  \Gamma' = \Gamma_{\mathrm{classic}}
  \;\Big[\;
    s_1 \mapsto \sub_{\prim{S}}^{0.47}{}_{\!0.025},\quad
    s_3 \mapsto \sub_{\prim{C}}^{0.53}{}_{\!0.003}
  \;\Big]
\]
or as a crossover product:
\[
  \Gamma' = \Gamma_{\mathrm{classic}}[{:}2] \mathbin{|} \Gamma_{\mathrm{void}}[2{:}]
\]
meaning: steps 1–2 from \emph{classic}, steps 3–4 from \emph{void}.

% ── 7. Fitness annotation ──────────────────────────────────────────────────
\section*{Fitness annotation}

A fully evaluated grammar is annotated:
\[
  \Gamma_{\mathrm{void}} \xrightarrow{\;\mathcal{F}\;}
  \bigl(\,f = 0.734,\; d_f = 2.51,\; t_{\min} = 0.82\,\mathrm{mm}\bigr)
\]
where $f$ is total fitness, $d_f$ the estimated fractal dimension,
and $t_{\min}$ the minimum wall thickness at $80\,\mathrm{mm}$ scale.

\end{document}
